% 火星（综述）
% license CCBYSA3
% type Wiki

本文根据 CC-BY-SA 协议转载翻译自维基百科\href{https://en.wikipedia.org/wiki/Mars}{相关文章}。

\begin{figure}[ht]
\centering
\includegraphics[width=6cm]{./figures/b86f152bd3571454.png}
\caption{火星的真实颜色\(^{\text{a}}\)，由“希望号”轨道器拍摄。画面中央可见塔尔西斯山脉，左侧为奥林匹斯山，右侧则是水手谷。} \label{fig_Mars_1}
\end{figure}
火星是距离太阳的第四颗行星。由于其呈橙红色的外观，它也被称为“红色星球”\(^{\text{22,23}}\)。火星是一颗类似沙漠的岩质行星，大气稀薄，主要成分为二氧化碳（CO\(_2\)）。在平均地表高度处，其大气压仅为地球的几千分之一，大气温度范围为 −153 至 20 °C（−243 至 68 °F）\(^{\text{24}}\)，宇宙辐射水平很高。火星保留了一部分水，存在于地下以及稀薄大气中，可形成卷云、雾气、霜冻、较大的极地永久冻土区与冰帽（伴随季节性的 CO\(_2\) 雪），但其表面不存在液态水体。其表面重力约为地球的三分之一，为月球的两倍。直径为 6,779 公里（4,212 英里），约为地球的一半，是月球的两倍，其表面积则相当于地球全部陆地面积的总和。

细微尘埃遍布火星表面和大气，在低重力条件下，即便是稀薄大气的微弱风力也能将其扬起并扩散开来。火星地形大致呈现显著的南北分界，即“火星二元性结构”：北半球主要是平坦、低洼的平原，而南半球则由遍布陨石坑的高地组成。从地质学上看，火星处于相对活跃的状态，地下会发生“火震”，但其上也分布着许多巨型而已灭绝的火山（最高的是奥林匹斯山，高 21.9 公里或 13.6 英里），以及太阳系中最大峡谷之一的水手谷（长约 4,000 公里或 2,500 英里）。火星有两颗天然卫星，形状小而不规则：火卫一（Phobos）与火卫二（Deimos）。由于其轴倾角为 25 度，火星像地球（轴倾角 23.5 度）一样经历季节变化。一个火星太阳年相当于 1.88 个地球年（687 个地球日），一个火星太阳日（sol）则等于 24.6 小时。

火星与其他行星一同形成于大约 45 亿年前。在火星的诺亚纪（Noachian period，约 45 亿年至 35 亿年前）期间，其地表经历了陨石撞击、峡谷形成、侵蚀作用、可能存在的海洋，以及磁层的消失。赫斯珀里亚纪（Hesperian period，自约 35 亿年前开始，结束于约 33–29 亿年前）以广泛的火山活动和导致巨大泄流河道形成的洪水为主要特征。亚马逊纪（Amazonian period）从其后延续至今，是当前地质过程的主要影响因素。由于火星的地质历史，对于火星过去或当前是否存在生命的探索仍是活跃的科学研究领域，其中一些可能的生命迹象仍需进一步确认。

由于在地球夜空中可被肉眼看到、呈红色并缓慢移动，火星自古以来就被人类观测，并在不同文化中具有多样的象征意义。1963 年，人类首次发射前往火星的探测器“火星 1 号”（Mars 1），但途中失联。第一项成功飞越火星的探测任务发生在 1965 年，由“水手 4 号”（Mariner 4）完成。1971 年，“水手 9 号”（Mariner 9）进入火星轨道，成为人类历史上首个进入除月球、太阳或地球以外天体轨道的航天器；同年，“火星 2 号”成为首个在火星发生非受控撞击的探测器，“火星 3 号”则实现了首次成功着陆。自 1997 年以来，探测器在火星上持续运行。有时，超过十个探测器会同时在火星轨道或地表开展活动，这一数量超过地球以外任意行星。

火星常被视为未来载人探索任务的目标，尽管目前尚无正式计划付诸实施。



